\chapter{Introduzione e contesto}
\label{chap:introduzione}

\section{Motivazione}
\label{sec:motivazione}

La gestione documentale aziendale rappresenta un caso paradigmatico in cui
l'affidabilit\`a operativa, la compliance normativa e l'efficienza dei processi
si intrecciano. In particolare, le imprese che operano in regime di fatturazione
elettronica (B2B/B2G) devono trattare flussi continuativi di documenti
fiscali, contrattuali e amministrativi in arrivo da molteplici canali, con
vincoli di tracciabilit\`a, integrit\`a, riservatezza e conservazione a lungo
termine. La gestione manuale di tali flussi si traduce sistematicamente in
ritardi di registrazione, errori di classificazione, smarrimento di allegati e
violazioni dei termini di pagamento.

L'introduzione di capacit\`a \emph{agentiche} basate su \emph{Large Language Models}
(LLM) consente di automatizzare il triage e la classificazione semantica dei
documenti, ma introduce nuove classi di rischio: allucinazioni, costi di
inferenza, perdita di accountability, azioni irreversibili eseguite con bassa
confidenza. Il progetto affronta questo trade-off costruendo un sistema in cui
l'agente \`e un \emph{operational assistant} vincolato da \emph{guardrail}
espliciti e sempre sottoposto, nei casi ambigui, ad approvazione umana.

\section{Obiettivi}
\label{sec:obiettivi}

Gli obiettivi del progetto, coerenti con i requisiti del corso, sono i seguenti:

\begin{itemize}[leftmargin=*]
    \item progettare un servizio distribuito \emph{scalabile} e \emph{affidabile}
          per la gestione di documenti aziendali;
    \item integrare uno \emph{strato operativo agentico}, basato su MCP, capace
          di assolvere a ruoli operativi distinti (interpretazione del contenuto,
          classificazione, escalation);
    \item definire una politica di sicurezza operativa esplicita che disciplini
          azioni automatiche, advisory e \emph{human-in-the-loop};
    \item dimostrare il comportamento del sistema in condizioni di fallimento
          controllato, con particolare riferimento a \emph{graceful degradation},
          \emph{budget compliance} e \emph{Recovery Time Objective}.
\end{itemize}

\section{Sintesi dei contributi}
\label{sec:contributi}

Il sistema realizzato integra:
\begin{enumerate}[leftmargin=*]
    \item un'architettura a microservizi containerizzata con \emph{S3-compatible object
          storage} (\textsc{SeaweedFS}) e \textsc{PostgreSQL} come piano dei metadati;
    \item un \emph{poller IMAP} con scheduling differenziato per utenti attivi e inattivi,
          idempotente sull'\textsc{UID} di messaggio;
    \item un \emph{agent worker} basato su \emph{Atomic Agents} e \emph{instructor},
          che produce \emph{structured outputs} e applica una soglia di confidenza per la
          filiazione automatica;
    \item un \emph{MCP server} che espone l'albero delle cartelle e l'inventario
          documentale al solo livello di \emph{read-only telemetry};
    \item una politica di sicurezza che instrada in \emph{inbox} umano ogni
          classificazione a bassa confidenza, garantendo reversibilit\`a totale delle
          decisioni dell'agente.
\end{enumerate}

\section{Struttura del documento}
\label{sec:struttura}

Il report \`e organizzato come segue. Il \Cref{chap:servizio} definisce il servizio,
gli utenti e gli scenari di guasto. Il \Cref{chap:architettura} descrive l'architettura
a microservizi e i flussi principali. Il \Cref{chap:stack} motiva le scelte tecnologiche.
Il \Cref{chap:database} dettaglia lo schema relazionale. Il \Cref{chap:api} documenta
le API REST esposte. Il \Cref{chap:agentico} approfondisce il layer agentico, il protocollo
MCP e la safety policy. Il \Cref{chap:sicurezza} tratta l'autenticazione, la crittografia
e l'isolamento multi-tenant. Il \Cref{chap:reliability} analizza i meccanismi di affidabilit\`a
e scalabilit\`a. Il \Cref{chap:observability} descrive logging, health check e audit trail.
Il \Cref{chap:devsecops} tratta la pipeline CI/CD e le pratiche DevSecOps. Il
\Cref{chap:testing} documenta la strategia di test. Il \Cref{chap:roi} presenta l'analisi
SLO, costi e Return on Agent. Il \Cref{chap:failure} analizza il comportamento del sistema
sotto condizioni degradate. Il \Cref{chap:conclusioni} riassume e prospetta gli sviluppi futuri.
